\documentclass{article}
\usepackage{amssymb, amsmath}
\usepackage[left=1.5cm, right=2cm]{geometry}
\begin{document}

\textbf{Public Key Infrastructure}

\textbf{Public Keys are very useful}

Secure web connections\\
Software signing\\
Secure messaging, email\\
Cryptocurrency, blockchains\\
How do we know the public key of an entity? And how can we trust that this entity is indeed who they claim to be and that they own a specific public key? Mainly: the key must be signed by a trusted Certificate Authority (CA)\\
Public Key Infrastructure (PKI) defines how to issue, manage, and use such certificates

\textbf{Public Key Certificates \& Authorities}

The big picture: when receiving a party's (the subject's) public key, it will be accompanied with a certificate. A valid cert means that the entity is who they claim to be and they own the corresponding public key\\
Certificate: signature by a CA over subject's public key and attributes\\
Attributes: identity (ID) and others (e.g. location)\begin{itemize}
	\item Validated by CA (liability?)
	\item Used by relying party for decisions
\end{itemize}
Root CA is most trusted in the hierarchy

\textbf{Certificates are all about Trust}

\textbf{Rogue Certs}

Rogue certificate: certificates that contain wrong or misleading information (so they \textit{should} fail PKI validation)\\
Attacker goals:\begin{itemize}
	\item Impersonate: website, phishing email, signed malware...
	\item Equivocating (same name): circumvent name-based security mechanisms such as blacklists, access-control..
\end{itemize}
Types of misleading names:\begin{itemize}
	\item Combo names: bank.com vs accts-bank.com, bank.accts.com
	\item Domain-name hacking: accts.bank.com vs accts-bank.com, accts.bank.co
	\item Homographic: paypal.com vs paypaI.com
	\item Typo-squatting: banc.com, baank.com
\end{itemize}

\textbf{PKI Failures}

Although the signature over the cert verifies correctly, there is still a failure and the cert may be revoked: this is called a PKI failure\\
PKI failures include:\begin{itemize}
	\item Corrupt CA
	\item Validation failure
	\item Exposed CA private key
	\item Cryptanalysis cert forgery: find collisions in HtS paradigm, or exploit some vulnerability in the digital signature scheme used for signing
\end{itemize}

\textbf{X.509 Certificates}

\textbf{The X.509 Standard Certificate Format}

Idea: signature binds public key to distinguished name (DN) and to other attributes\\
Used widely despite complaints about its complexity: SSL/TLS, code-signing, IP-Sec, ...\\
Signed Fields:\begin{enumerate}
	\item Version
	\item Cert serial number
	\item Signature Algorithm Object Identifier (OID)
	\item Issuer Distinguished Name (DN)
	\item Validity period (make sure it's time-limited)
	\item Subject (user) DN
	\item Subject public key info
\end{enumerate}
Then, signature

\textbf{X.509 Certs and Subject Identifiers}

V1: Distinguished Name (for subject and issuer)
V2: Unique Identifiers (for subject and issuer)
V3: Extensions (used in practice)

\textbf{Distinguished Names}

Most certs contain identifiers\\
Influenced by telecom providers: phone directory services are based on common names\\
Basic goals of identifiers:\begin{itemize}
	\item Meaningful (to humans)
	\item Unique identification of entity (owner)
	\item Decentralized, with accountability
\end{itemize}

\textbf{The Identifiers Trilemma}

Achieving the three goals: Meaningful, Unique, Decentralized, seems very challenging\\
Examples of achieving any \underline{two} of the goals:\begin{itemize}
	\item Unique + Meaningful: URL, email
	\item Meaningful + Decentralized: common name
	\item Unique + Decentralized: hash of key
\end{itemize}

\textbf{X.509 Validation of Certificate Paths}

Have Root CA, with ICAs (intermediary CA) and LCAs (leaf CAs)\\
Only leaf CAs interact with end users.\\
Simply validate all certs in the chain all the way to the root CA\\
The root CA (self-signed) cert is in the root store in Alice's browser

\textbf{Certificate Revocation}

Reasons for revoking:\begin{itemize}
	\item Security issues: key compromised, CA compromised
	\item Administrative issues: change name, location
\end{itemize}
Techniques

\textbf{Certificate Revocation List (CRL)}
A CRL is simply a list of revoked certs, distributed periodically by CAs\\
If CRLs contain all revoked certs (which did not expire) it could be huge!\\
CRLs are not immediate: who is responsible until CRL is distributed. Frequent CRLs $\rightarrow$ even more overhead

\textbf{CRL Optimization Solutions}

CRL distribution point: split certs to several CRLs\\
Authorities Revocation List (ARL): list only revoked CAs (very coarse grain, revoking CAs as a whole but not necessarily individuals)\\
Delta CRL: only new revocations since last "base" CRL-- need to keep CRLs for a long period to check deltas\\
Another more efficient approach is the Online Certificate Status Protocol (OCSP)

\textbf{Online Certificate Status Protocol (OCSP)}

Improve efficiency and freshness compared to CRLs\\
Client asks CA about cert during handshake\\
CA signs response (real-time)\\
Adds additional delay to access time\\
Tradeoff: less cost, but higher latency

\textbf{OCSP Challenges}

Privacy (expose domain and client to CA), load on CA, response delay, reliability (what if CA fails)\\
Ambiguity: when an OCSP server (or CA) cannot resolve the request, it replies with "certificate status is unknown"\\
Reliability or failed requests: client failed to establish connection with OCSP server, or client's request is invalid (not signed, or not authorized) so no response will be received.

\textbf{Ambiguous/Failed OCSP Responses}

What should client do?\\
Wait forever? Unrealistic\\
Hard-fail: terminate connection since certificate is unknown/not received (safe!)\\
Ask user: application displaying message asking user how they want to proceed\\
Soft-fail: pretend that a response has been received and continue as if the cert is not revoked (common choice for browsers, but a MitM attacker may block the OCSP response to make a revoked cert go through)

\textbf{Conclusion}

PKI is an essential component of the internet\\
Yet, it is a complicated module with many issues related to security, privacy, and performance\begin{itemize}
	\item To many, this is a solved problem (but that is not the case)
	\item Several open questions related on how to detect rogue certs, handle CA failure, revocation ,etc, how to audit these parties, reduce trust..
	\item How to handle all these issues efficiently? We all want a highly responsive Web!
\end{itemize}


\end{document}
